\origpage{141}{GRAPH MINORS. XV. GIANT STEPS}
\noindent $\theta > 6\pi + 19$ and $H, \mathscr{T}$ is a $\Sigma$-span, it follows that $(A^* \cap H, B^* \cap H) \in \mathscr{T}$. Consequently $(A \cap H, B \cap H) \notin \mathscr{T}$, because $(A^* \cap H) \cup (A \cap H) \supseteq (A^* \cap H) \cup (B^* \cap H) = H$, and so $(A, \Omega) \notin \mathscr{S}_0$, because $\mathscr{S}_0$ is $\mathscr{T}$-local. Thus there exists $i$ with $1 \le i \le k$ such that $(A, \Omega) \in \mathscr{S}_i$. Then $B^* \subseteq A \subseteq G_i^*$. However, by [5, Theorem (7.3)], there exists $e \in E(B^* \cap H)$ with
\begin{equation*}
d(r_i, e) \ge \theta - |V((A^* \cap H) \cap (B^* \cap H))| \ge \theta - |V(A^* \cap B^*)| \ge \theta - 6\pi - 19.
\end{equation*}
Since $e \in E(H)$, there exists a unique $(A', \Omega') \in \mathscr{S}_0$ with $e \in E(A')$; and since $e \in E(G_i^*)$ it follows that $\alpha_0(A', \Omega') \subseteq \Delta_i$. Since $\alpha_0$ is compatible with $H$, we deduce that $e \subseteq \Delta_i$, and so $d(r_i, e) \le 3\beta + 6$ by (7.5)(ii). But $3\beta + 6 < \theta - 6\pi - 19$, a contradiction. This proves (4).

Define $\alpha(A, \Omega)$ for $(A, \Omega) \in \mathscr{S}$ as follows: if $(A, \Omega) \in \mathscr{S}_0$, let $\alpha(A, \Omega) = \alpha_0(A, \Omega)$; and otherwise, choose $i$ with $1 \le i \le k$ such that $(A, \Omega) \in \mathscr{S}_i$, and let $\alpha(A, \Omega) = \alpha_i(A, \Omega)$. For $v \in V(\mathscr{S})$, define $\alpha(v)$ by: if $v \in \bar{\Omega}$ for some $(A, \Omega) \in \mathscr{S}_0$, let $\alpha(v) = \alpha_0(v)$; and otherwise, choose $i$ with $1 \le i \le k$ such that $v \in \bar{\Omega}$ for some $(A, \Omega) \in \mathscr{S}_i$, and let $\alpha(v) = \alpha_i(v)$.

\hspace*{12pt}(5) \textit{For $1 \le i \le k$, if $v \in V(\mathscr{S}_i)$ then $\alpha(v) = \alpha_i(v)$.}

\textit{Subproof.} Suppose first that $v \in \bar{\Omega}$ for some $(A, \Omega) \in \mathscr{S}_0$. Then $\alpha(v) = \alpha_0(v)$; but by (1), $v \in \bar{\Omega}_i^*$ and so $\alpha_i(v) = \alpha_0(v) = \alpha(v)$ as required. We assume then that there is no such $(A, \Omega)$. Choose $1 \le j \le k$ and $(A, \Omega) \in \mathscr{S}_j$ with $v \in \bar{\Omega}$, such that $\alpha(v) = \alpha_j(v)$. Then $v \in V(G_i^* \cap G_j^*)$, and so $i = j$ by (7.6). Hence $\alpha_i(v) = \alpha(v)$. This proves (5).

\hspace*{12pt}(6) \textit{$\alpha \in bd(\alpha(A, \Omega))$ for each $(A, \Omega) \in \mathscr{S}$ and each $v \in \bar{\Omega}$.}

\textit{Subproof.} If $(A, \Omega) \in \mathscr{S}_0$, then $\alpha(A, \Omega) = \alpha_0(A, \Omega)$ and $\alpha(v) = \alpha_0(v) \in bd(\alpha_0(A, \Omega))$, as required. We assume then that $(A, \Omega) \in \mathscr{S}_i$ for some $i \ge 1$. By (5), it follows that $\alpha(v) = \alpha_i(v) \in bd(\alpha_i(A, \Omega))$. This proves (6).

\hspace*{12pt}(7) \textit{For distinct $(A, \Omega), (A', \Omega') \in \mathscr{S}$, if $x \in \alpha(A, \Omega) \cap \alpha(A', \Omega')$ then $x = \alpha(v)$ for some $v \in \bar{\Omega} \cap \bar{\Omega}'$.}

\textit{Subproof.} If $(A, \Omega), (A', \Omega') \in \mathscr{S}_0$, then $x \in \alpha_0(A, \Omega) \cap \alpha_0(A', \Omega')$, and so $x = \alpha_0(v) = \alpha(v)$ for some $v \in \bar{\Omega} \cap \bar{\Omega}'$. We may assume then that $(A, \Omega) \notin \mathscr{S}_0$. Choose $i$ with $1 \le i \le k$ such that $(A, \Omega) \in \mathscr{S}_i$ and $\alpha(A, \Omega) = \alpha_i(A, \Omega) \subseteq \Delta_i$. Suppose that $(A', \Omega') \notin \mathscr{S}_0$. Then there exists $j$ with $1 \le j \le k$ such that $(A', \Omega') \in \mathscr{S}_j$ and $\alpha(A', \Omega') = \alpha_j(A', \Omega')$; and since $\alpha_j(A', \Omega') \subseteq \Delta_j$ and $\alpha_i(A, \Omega) \subseteq \Delta_i$, and so $x \in \Delta_i \cap \Delta_j$, it follows that $i = j$. Hence there exists $v \in \bar{\Omega} \cap \bar{\Omega}'$ with $\alpha_i(v) = x$, and so $\alpha(v) = x$, by (5). We may therefore assume that $(A', \Omega') \in \mathscr{S}_0$. Then $\alpha(A, \Omega) \subseteq \Delta_i$, and $\alpha(A', \Omega') \cap \Delta_i \subseteq bd(\Delta_i)$, and so $x \in bd(\Delta_i)$. Since $bd(\Delta_i)$ is $\alpha_0$-normal and $x \in \alpha_0(A', \Omega')$, it follows that $x = \alpha_0(v)$ for some $v \in \bar{\Omega}'$; and hence $v \in \bar{\Omega}_i^*$. Consequently $v \in V(\mathscr{S}_i)$, and $\alpha_i(v) = \alpha_0(v) = x$. Since $x \in \alpha_i(A, \Omega)$ we deduce
